\documentclass[10pt,a4paper]{amsart}
\usepackage[margin=19mm]{geometry}
\usepackage{amsmath,amssymb,mathtools}
\usepackage[scr=rsfs,cal=euler]{mathalfa}
\usepackage{xfrac}
\usepackage[unicode,hidelinks,bookmarks=false]{hyperref}
\newcommand{\ie}{i.e.\ }
\newcommand{\cf}{cf.\ }
\newcommand{\resp}{respectively\ }
\newcommand{\Subsection}{\subsection}
\providecommand{\nobreakdash}{\nobreak\hbox{-}}
\setlength{\emergencystretch}{2em}
\multlinegap0pt
\usepackage{dsfont}
\usepackage{amssymb}
\usepackage{bm}
\usepackage{tikz}
\usepackage{mathtools}
\newtheorem{lemma}{Lemma}
\newcommand{\R}{\mathbb{R}}
\newcommand{\ep}{\epsilon}
\newcommand{\p}{\mathcal{P}}
\title{{\bf Uniqueness of positive solutions of double-power nonlinear stationary Schr\"odinger equations:} \\the fixed frequency case and the fixed mass case}
\author{Linjie Song, Wenming Zou}
\date{Source: arXiv:2609.15882v1 (CC BY 4.0)}
\begin{document}
\maketitle

\noindent\emph{Source and licence.} {\bf Uniqueness of positive solutions of double-power nonlinear stationary Schr\"odinger equations:} \\the fixed frequency case and the fixed mass case. Version arXiv:2609.15882v1 (CC BY 4.0), distributed by arXiv under the Creative Commons Attribution 4.0 International licence, http://creativecommons.org/licenses/by/4.0/, as recorded in the arXiv licence metadata for this exact version and shown on the abstract page https://arxiv.org/abs/2609.15882v1. Full licence notice: \texttt{licenses/arxiv-2609.15882-CC-BY-4.0.txt}.\par\smallskip

\noindent\textit{Editorial scope.} Counted unit: the article's lemma lemestimate (source0.tex lines 433--439). The article's complete original proof of it is delivered with it (source0.tex lines 441--640, one \texttt{proof} environment of 200 lines). No mathematics is written, repaired or re-derived: the statement and the proof are the article's own lines, in the article's own notation, and no retained line is substituted or re-flowed.  Retained uncounted as the source-local context the proof needs: lines 371--373.  Nothing was omitted from any retained span, so the package carries no omission record.  No retained line contains a citation, so the wrapper carries no bibliography and no bibliography entry.  No retained line contains a figure environment, an image-inclusion command or drawing code, so no image is delivered and the package declares no asset dependency.  Editorial formatting only, in addition: an AMS \texttt{amsart} preamble that restates the article's own declarations and macros used by the retained lines, namely the environments \texttt{lemma} on separate counters, together with the standard packages those lines need.  The retained environments print in this document as Lemma 1 in order of appearance, whereas the article prints the same items with the numbering of its own full document.  The complete native source of the article as distributed in the arXiv e-print for this exact version is bundled unchanged as \texttt{sources/210460/source0.tex}.
	\begin{align} \label{eqofv}
		-\Delta v + \eta v = (1-\theta)v^p + \theta v^q \quad \text{ in } \R^N,
	\end{align}

	\begin{lemma} \label{lemestimate}
		Suppose $N \ge 2$ and $1 < q < p < 2^* -1$. Then
		\begin{align} \label{usefulestimate1}
			y_r > 0, \quad h_r < 0, \quad J - h < k_r < J
		\end{align}
		for $r > 0$.
	\end{lemma}

	\begin{proof}
		Let
		\begin{align*}
			\p_0 = r^N\left(v_r^2 + 2\mathcal{F}(v)\right) + (N-2)r^{N-1}vv_r
		\end{align*}
		where
		\begin{align*}
			\mathcal{F}(t) = -\frac{\eta}{2}t^2 + \frac{1-\theta}{p+1}t^{p+1} + \frac{\theta}{q+1}t^{q+1}.
		\end{align*}
		From direct computations we derive that
		\begin{align*}
			\p_0'(r) = r^{N-1}v^2\left(-2\eta + (1-\theta)\frac{N+2-p(N-2)}{p+1}v^{p-1} + \theta\frac{N+2-q(N-2)}{q+1}v^{q-1}\right).
		\end{align*}
		Since $v(r)$ is decreasing in $r > 0$ and $q < p < 2^* -1$, the function
		\begin{align*}
			-2\eta + (1-\theta)\frac{N+2-p(N-2)}{p+1}v(r)^{p-1} + \theta\frac{N+2-q(N-2)}{q+1}v(r)^{q-1}
		\end{align*}
		is decreasing in $r > 0$ and has at most one zero point. Furthermore, using $\p_0(0) = 0$ and $\p_0(r) \to 0$ as $r \to \infty$, we obtain $\p_0(r) > 0$ for all $r > 0$. On the other hand,
		\begin{align*}
			\p_0 = r^{N-2}v^2\left(ry_r - r^2J\right).
		\end{align*}
		Therefore,
		\begin{align*}
			y_r > rJ > 0.
		\end{align*}

		Next we prove $h_r < 0$. Since $v(0) = 1$, we have
		\begin{align*}
			h(0) = - \lim_{r \to 0}\frac{v_r(r)}{r} = -v_{rr}(0) =: h_0 > 0.
		\end{align*}
		Using \eqref{eqofv} we get
		\begin{align*}
			-N v_{rr}(0) = (1-\theta)v(0)^{p} + \theta v(0)^{q} - \eta v(0) = 1 - \eta.
		\end{align*}
		Thus $h_0 = (1-\eta)/N$. By the regularity and radial symmetry of $v(x)$, we can write
		\begin{align*}
			h(r) = h_0 + \mu r^2 + o(r^2), \quad h_r(r) = 2\mu r + o(r).
		\end{align*}
		Using $\xi_r = rh$, we get
		\begin{align*}
			\xi(r) = \int_0^r sh(s)ds = \int_0^r s\left(h_0 + \mu s^2 + o(s^2)\right)ds = \frac{h_0}{2}r^2 + O(r^4).
		\end{align*}
		Moreover, $F_\xi(0) = -P(0)$. Thus
		\begin{align*}
			F(\xi) = F(0) + F_\xi(0)\xi + o(\xi) = 1 - P(0)\xi  + o(\xi) = 1 - \frac{h_0P(0)}{2}r^2 + o(r^2).
		\end{align*}
		By $v_r/v = -rh$, we have
		\begin{align*}
			\frac{v_{rr}}{v} = -h - rh_r + r^2h^2.
		\end{align*}
		On the other hand,
		\begin{align*}
			-v_{rr} - \frac{N-1}{r}v_r + \eta v = F(\xi)v.
		\end{align*}
		Therefore, we can derive that
		\begin{align} \label{eqofh_r}
			rh_r = r^2h^2 + F -\eta - Nh.
		\end{align}
		Note that
		\begin{align*}
			& rh_r = r(2\mu r + o(r)) = 2\mu r^2 + o(r^2), \\
			& r^2h^2 = r^2(h_0 + \mu r^2 + o(r^2))^2 = h_0^2 r^2 + o(r^2),
		\end{align*}
		and
		\begin{align*}
			F - \eta - Nh & = 1 - \frac{h_0P(0)}{2}r^2 + o(r^2) - \eta - N (h_0 + \mu r^2 + o(r^2)) \\
			& = 1 - \eta - Nh_0 - \frac{h_0P(0) + 2N\mu}{2}r^2 + o(r^2) \\
			& = - \frac{h_0P(0) + 2N\mu}{2}r^2 + o(r^2).
		\end{align*}
		Comparing the coefficient of $r^2$, one can see that
		\begin{align*}
			2\mu = h_0^2 - \frac{h_0P(0) + 2N\mu}{2},
		\end{align*}
		namely
		\begin{align*}
			\mu = \frac{h_0(2h_0 - P(0))}{2(N+2)}.
		\end{align*}
		Thus
		\begin{align} \label{eqextending}
			h(r) = h_0 + \frac{h_0(2h_0 - P(0))}{2(N+2)} r^2 + o(r^2), \quad h_r(r) = \frac{h_0(2h_0 - P(0))}{N+2} r + o(r).
		\end{align}

		We claim that $2h_0 \le P(0)$. If it is not in this case, when $r > 0$ is small enough we have $h'(r) > 0$. On the other hand, as $r \to \infty$,
		\begin{align*}
			0 < h(r) = -\frac{v_r(r)}{rv(r)} \le \frac{C}{r} \to 0.
		\end{align*}
		This yields the existence of $R > 0$ such that
		\begin{align*}
			h_r(r) > 0 ~~ (0 < r < R), \quad\quad h_r(R) = 0.
		\end{align*}
		We also have $h_{rr}(R) \le 0$. Now we introduce
		\begin{align*}
			Q(r) := 2h(r) - P(\xi(r)).
		\end{align*}
		Direct computations show that
		\begin{align*}
			Q_r = 2h_r - P_\xi \xi_r = 2h_r - rh P_\xi.
		\end{align*}
		Since $h > 0$ and $P_\xi < 0$, we get $Q_r > 0$ in $r \in (0,R)$ and so $Q(R) > Q(0) = 2h_0 - P(0) > 0$. Differentiating \eqref{eqofh_r} with respect to $r$, we deduce that
		\begin{align*}
			rh_{rr} + h_r = 2r^2hh_r + 2rh^2 + F_\xi \xi_r - Nh_r = 2r^2hh_r + 2rh^2 - rhP - Nh_r,
		\end{align*}
		namely
		\begin{align} \label{eqofhrr}
			h_{rr} = hQ + \left(2rh - \frac{N+1}{r}\right)h_r.
		\end{align}
		This shows $h_{rr}(R) = h(R)Q(R) > 0$, a contradiction! Thus the claim holds true.

		In the following we aim to prove $h_r \le 0$ for $r > 0$. Suppose by contradiction that the set
		\begin{align*}
			E = \big\{r > 0: h_r(r) > 0\big\}
		\end{align*}
		is not empty. Take a connected component $(\alpha,\beta) \subset E$. It is clear that $h_r(r) > 0$ for $\alpha < r < \beta$. Since $h > 0$ and $h(r) \to 0$ as $r \to \infty$, we have $\beta < \infty$. At the right endpoint $r = \beta$,
		\begin{align*}
			h_r(\beta) = 0, \quad\quad h_{rr}(\beta) \le 0.
		\end{align*}
		If $\alpha > 0$, at the left endpoint $r = \alpha$ we also have
		\begin{align*}
			h_r(\alpha) = 0, \quad\quad h_{rr}(\alpha) \ge 0.
		\end{align*}
		Using \eqref{eqofhrr} we get
		\begin{align*}
			h_{rr}(\alpha) = h(\alpha)Q(\alpha).
		\end{align*}
		Together with $h(\alpha) > 0$ we obtain $Q(\alpha) \ge 0$. Reasoning as before we have
		\begin{align*}
			Q(\beta) > Q(\alpha) \ge 0.
		\end{align*}
		Using \eqref{eqofhrr} again,
		\begin{align*}
			h_{rr}(\beta) = h(\beta)Q(\beta) > 0,
		\end{align*}
		contradicting $h_{rr}(\beta) \le 0$. If $\alpha = 0$, we have proved that $Q(0) = 2h_0 - P(0) \le 0$. If $Q(0) < 0$, by \eqref{eqextending} one gets $h_r(r) < 0$ near $r = 0$, implying that $(0,\beta) \not\subset E$. Therefore, $Q(0) = 0$ and it still holds that
		\begin{align*}
			Q(\beta) > Q(0) = 0
		\end{align*}
		and
		\begin{align*}
			h_{rr}(\beta) = h(\beta)Q(\beta) > 0
		\end{align*}
		a contradiction! Thus the set $E$ must be empty, namely $h_r \le 0$ in $r > 0$.

		Further, we will eliminate the possibility that $h_r$ has a zero point. Arguing by contradiction, we assume that $h_r(R) = 0$ for some $R > 0$. Together with $h_r \le 0$ we get $h_{rr}(R) = 0$. By \eqref{eqofhrr} and using $h(R) > 0$ we have $Q(R) = 0$. At $r = R$,
		\begin{align*}
			Q_r(R) = 2h_r(R) - Rh(R)P_\xi(\xi(R)) = - Rh(R)P_\xi(\xi(R)) > 0.
		\end{align*}
		Differentiating \eqref{eqofhrr} with respect to $r$, we deduce
		\begin{align*}
			h_{rrr} = h_rQ + hQ_r + A_r(r)h_r + A(r)h_{rr}
		\end{align*}
		where
		\begin{align*}
			A(r) = 2rh - \frac{N+1}{r}.
		\end{align*}
		At $r = R$,
		\begin{align*}
			h_{rrr}(R) = h(R)Q_r(R) > 0.
		\end{align*}
		Expanding $h_r$ at $r = R$ we have
		\begin{align*}
			h_r(R + \ep) = h_r(R) + h_{rr}(R)\ep + \frac12 h_{rrr}(R)\ep^2 + o(\ep^2) = \frac12 h_{rrr}(R)\ep^2 + o(\ep^2).
		\end{align*}
		For small $\ep > 0$ we derive that $h_r(R + \ep) > 0$, contradicting $h_r \le 0$. This contradiction yields that $h_r < 0$ for all $r > 0$.

		Finally, by
		\begin{align*}
			y_r = (rk)_r = r k_r + k = rk_r + rh > rJ
		\end{align*}
		we get $k_r > J - h$. It remains to prove $k_r < J$. Let
		\begin{align*}
			\p_* = r^N\left(v_r^2 + 2\mathcal{F}(v)\right) + (N-1)r^{N-1}vv_r.
		\end{align*}
	    Then we have
		\begin{align*}
			\p_* = r^Nv^2 \left(k_r - J\right), \quad \p_*'(r) = r^{N-1}v^2 \Sigma
		\end{align*}
		where
		\begin{align*}
			\Sigma(r) = k^2 - \eta + (1-\theta)\frac{N+1-p(N-1)}{p+1}v^{p-1} + \theta\frac{N+1-q(N-1)}{q+1}v^{q-1}.
		\end{align*}
		We first prove $\p_* \le 0$ in $r > 0$. Otherwise, since $\p_*(0) = 0$ and $\p_*(r) \to 0$ as $r \to \infty$, we can assume the existence of $R > 0$ such that
		\begin{align*}
			\p_*(R) > 0, \quad\quad \p_*'(R) = 0.
		\end{align*}
		Clearly, $\Sigma(R)  = 0$. Since
		\begin{align*}
			k_r = k^2 - \frac{v_{rr}}{v} = k^2 - \eta + (1-\theta)v^{p-1} + \theta v^{q-1} - (N-1)h,
		\end{align*}
		at $r = R$ we get
		\begin{align*}
			k_r(R) = N J(\xi(R)) - (N-1)h(R), \quad Rh_r(R) = k_r(R) - h(R) = N\left(J(\xi(R)) - h(R)\right).
		\end{align*}
		By $h_r < 0$ we get $h(R) > J(\xi(R))$. Thus
		\begin{align*}
			\p_*(R) = (N-1)R^{N}v(R)^2\left(J(\xi(R)) - h(R)\right) < 0,
		\end{align*}
		a contradiction!

		Now we suppose by contradiction that there exists $r_0 > 0$ such that $\p_*(r_0) = 0$. In view of $\p_* \le 0$, we know $\p_*'(r_0) = 0$ immediately. Arguing as before, we get $\p_*(r_0) < 0$ and find a contradiction. This shows $\p_* < 0$, namely $k_r < J$ in $r > 0$. The proof is complete.
	\end{proof}

\end{document}
